\section{モジュラー型 \texorpdfstring{$(2,2)$}{(2,2)} の全層の帰属}

本節では
\[
 J=\partial_x\wedge\partial_y+\partial_z\wedge\partial_w,
 \qquad
 \mathcal Q=\{Q\text{：斉次二次二重ベクトル};\ [J,Q]=[Q,Q]=0\}
\]
を考える。互換条件 $[J,Q]=0$ の解空間は45次元である。
以下の接空間はすべてこの45変数空間で計算する。
固定トレースや固定 Jordan 型の接空間だけを用いた成分判定ではない。

v6で、トレース $s=Q^{12}+Q^{34}=xz$ を固定した被約ファイバーは
17個の4次元既約成分に分解され、中心化群の成分群を考慮した
$\operatorname{Sp}_4$ 飽和は11個の10次元既約層となった。
そのうち開二枝の飽和は、既に $\mathcal Q$ の成分と証明されている。
今回、残る10層の\emph{全 $\mathcal Q$ 内での帰属}を決定した。

\subsection{共通正規形と結論の一覧}

閉15枝の共通正規形は、$p=p_3x^2+p_4xz+p_5z^2$ とおくと
\begin{align*}
 Q^{12}&=(1-d)xz-vx^2-uz^2,&Q^{13}&=0,\\
 Q^{14}&=az^2-4uxz+\frac{c-3d+2}{2}x^2,
 &Q^{23}&=bx^2-4vxz-\frac{c+d}{2}z^2,\\
 Q^{24}&=p+(2d-c-1)xy+czw+2uyz+2vxw,
 &Q^{34}&=dxz+vx^2+uz^2.
\end{align*}
各層の代表点では $(p_3,p_4,p_5)=(1,2,3)$ とする。
$C_\pm$ では
\[
 u=\frac1{16},\quad v=-\frac1{16},\quad
 c=\pm\frac18,\quad d=\frac12,\quad a=-\frac c2,\quad b=\frac c2,
\]
$D_+$ では
\[
 u=\frac16,\quad v=-\frac1{36},\quad c=\frac{\sqrt3}{18},\quad
 d=\frac12+c,\quad a=-3c,\quad b=-\frac c{12}
\]
をとる。他の枝の記号はv6と同じである。

\begin{center}
\begin{tabular}{c|c|c|c|c}
層（中心化群による同一視）&次元&$\dim T\mathcal Q$&全 $\mathcal Q$ での帰属&一般横断環\\\hline
開二枝&10&10&v6の成分&$K$\\
$C_+$&10&11&新成分&$K[\varepsilon]/(\varepsilon^4)$\\
$C_-$&10&12&新成分&$K[u,v]/(u,v)^2$\\
$D_\pm$&10&11&新成分&$K[\varepsilon]/(\varepsilon^3)$\\
$U_0,V_1$&10&11&$G_2$ の境界&---\\
$U_1,V_0$&10&12&新成分&$K[u,v]/(u,v)^2$\\
$L_0$&10&11&$G_2$ の境界&---\\
$L_1$&10&11&$G_1$ の境界&---\\
$A_0,B_0$&10&11&新成分&$K[\varepsilon]/(\varepsilon^3)$\\
$A_1,B_1$&10&11&$G_1$ の境界&---\\
$AB$&10&12&新成分&$K[u,v]/(u,v)^2$
\end{tabular}
\end{center}
ここで $K$ は各成分の関数体であり、横断環は一般点の形式近傍での主張である。
したがって $(2,2)$ 型から得られる全成分は7個、そのうち今回の追加は6個である。
表の横断環は、単に接空間次元を数えた推測ではなく、以下の形式方程式の計算による。

\subsection{二方向が二次で止まる三成分}

$C_-,U_1,AB$ の各代表コアを固定し、$p_3,p_4,p_5$ を変数とする。
\texttt{certify\_closed\_squarezero.py} は有理関数体
$\mathbb Q(p_3,p_4,p_5)$ 上で、全 Jacobi 微分の階数が33、
実際の族と $\operatorname{Sp}_4$ 作用の微分の階数が10であることを計算する。
証拠点でも同じ階数であるため、近傍で接空間の核と実際の族に対する商はベクトル束となる。
商の階数は2である。

その基底を $R,S$ とすると、全 Jacobi 微分の余核における
\[
 \operatorname{Jac}(R),\qquad
 \operatorname{polar}(R,S),\qquad
 \operatorname{Jac}(S)
\]
は独立である。保存した3次小行列の行列式は、順に
\[
 \frac{526848}{25},\qquad -184320,\qquad 98304
\]
である。ここで $\operatorname{polar}(R,S)=
\operatorname{Jac}(R+S)-\operatorname{Jac}(R)-\operatorname{Jac}(S)$ とする。

実際の滑らかな10次元族を形式座標で平坦化し、独立な線形方程式を消去する。
普遍一次接ベクトルが存在することから、残るイデアル $I$ は
$\mathbb C[[t_1,\ldots,t_{10},u,v]]$ 内で $(u,v)^2$ に含まれる。
二次障害の上記3次行列が可逆なので、逆の包含も成立し、
\[
 I=(u,v)^2
\]
となる。これにより全成分性、一般被約台の滑らかさ、一般厚さ3が同時に従う。
$C_-,AB$ では $\operatorname{Sp}_4\times\mathbb G_m$ 軌道の次元は9であるため、
群軌道だけで普遍化したとは主張しない。
有理関数体上の階数計算が、残りのモジュライ方向も含む族上の議論を保証している。

\subsection{一方向の形式スライスと高次障害}

ここで用いる形式スライスの条件を明記する。
滑らかな群 $G$ が滑らかな空間 $V$ の部分スキーム $X$ を保ち、
$q\in X$ の実際の群軌道の次元を $d$、全接空間の次元を $d+1$ とする。
軌道接空間に対する余剰方向 $R$ と、全接空間の補空間 $W\subset T_qV$ を選ぶ。
群の安定化群に横断する局所切断と線形スライス $q+KR+W$ から $V$ への作用写像は、
微分が同型なので形式近傍で同型となる。
したがって、実際の群軌道方向とスライス上の方程式を分離できる。

$W$ 上の微分が可逆となる独立方程式を選ぶと、形式陰関数定理により
\[
 w(\varepsilon)=\varepsilon^2w_2+\varepsilon^3w_3+\cdots
\]
が一意に定まる。残る\emph{すべての}方程式にこれを代入し、
$k-1$ 次までの係数が零、$k$ 次に非零係数があれば、
一変数形式環の残余イデアルは正確に $(\varepsilon^k)$ である。
任意に選んだ二次補正だけを調べるのではなく、線形スライスと陰関数解の一意性によって
高次補正の曖昧さを除いている。

$C_+,D_+$ では、全 Jacobi 階数34、$\operatorname{Sp}_4$ 軌道次元9、
$\operatorname{Sp}_4\times\mathbb G_m$ 軌道次元10であり、上記の条件を満たす。
\texttt{one\_normal\_higher\_jets.py} は45変数から34次元の $W$ を選び、
80個の全アフィン Jacobi 係数を各次数で保持する。
$C_+$ では二次・三次の全残余が零で、四次に初めて非零となる。
その最初の保存係数は $164560896$ である。
$D_+$ では二次残余が零で、三次に初めて非零となる。
従って一般横断環は、それぞれ
\[
 K[\varepsilon]/(\varepsilon^4),\qquad
 K[\varepsilon]/(\varepsilon^3)
\]
である。$D_+$ の全計算は正確な代数体 $\mathbb Q(\sqrt3)$ 上で行い、
未整理の根号式をそのまま非零と判定しない。

\subsection{残る四層の実際の多項式変形}

既知の11次元成分 $G_1,G_2$ の正規形を、成分の順序 $(12,13,14,23,24,34)$ で
\begin{align*}
 N_1(A,D,E)&=(Axy,Exz,-Exw,-Eyz,Eyw,Dzw),\\
 N_2(A,B,E)&=(Axy+Bzw,Exz,-Exw,-(B+E)yz,(B+E)yw,Bzw)
\end{align*}
とする。これらは任意の係数で Jacobi と $J$ 互換性を満たす。

対称行列 $K_0=\left(\begin{smallmatrix}a_0&b_0\\b_0&d_0\end{smallmatrix}\right)$ と
$L=\left(\begin{smallmatrix}\ell&m\\m&n\end{smallmatrix}\right)$ を用い、
\[
 S(t):(x,y,z,w)\longmapsto
 \left(x,y+(a_0/t+\ell)x+(b_0/t+m)z,
 z,w+(b_0/t+m)x+(d_0/t+n)z\right)
\]
とおく。$t\ne0$ でこれはシンプレクティックである。
また
\[
 A=t\alpha+t^2a_2,\qquad D\text{ または }B=t\beta+t^2b_2,
 \qquad E=t\eta+t^2e_2
\]
とする。以下の選択では $S(t)$ による変換後の負冪がすべて消え、
$Q(t)$ は $t$ の多項式になる。
\begin{center}
\begin{tabular}{c|c|c|c|c}
極限&元の成分&$(\alpha,\beta,\eta)$&$(a_0,b_0,d_0)$&追加座標変換\\\hline
$L_1$&$G_1$&$(1,1,e)$&$(0,1/2,0)$&恒等\\
$L_0$&$G_2$&$(-1,1,e)$&$(0,1,0)$&恒等\\
$A_1$&$G_1$&$(1,1,1/2)$&$(0,1/2,1)$&$T_A$\\
$U_0$&$G_2$&$(1,1,-1)$&$(1,1,1)$&$T_U$
\end{tabular}
\end{center}
ここで
\[
 T_A(x,y,z,w)=(x+z,y,z,w-y),\qquad
 T_U(x,y,z,w)=(x+2z,-y+w,x+z,2y-w).
\]
具体的には
\[
 Q(t)(v)=T S(t)^{-1}N_i(t)\bigl(S(t)T^{-1}v\bigr)S(t)^{-\mathsf T}T^{\mathsf T}
\]
である。

六つの補助係数 $(\ell,m,n,a_2,b_2,e_2)$ は、極限の $Q^{24}$ に任意の
$p_3x^2+p_4xz+p_5z^2$ を付加できる。
対応する線形写像の3次小行列は、表の順に
\[
 -\frac{(2e-1)(2e+1)}{16},\quad -4e(e+1),\quad -\frac18,\quad4
\]
である。従って一般パラメータで全磁気項を実現する。
$L_0,L_1$ の例外的な有限個の $e$ は閉包をとって含まれる。
$A_1,U_0$ の非零コア・パラメータはシンプレクティックな対角作用で動かせる。
これにより
\[
 \overline{\operatorname{Sp}_4L_1},\quad \overline{\operatorname{Sp}_4A_1}\subset G_1,
 \qquad
 \overline{\operatorname{Sp}_4L_0},\ \overline{\operatorname{Sp}_4U_0}\subset G_2
\]
が全被約閉包について成立する。
\texttt{certify\_boundary\_deformations.py} は多項式性、全 Jacobi 恒等式、互換性、
三つの磁気係数への全射性をすべて記号計算で検証する。

\subsection{\texorpdfstring{$A_0$}{A0} から得られる全 \texorpdfstring{$\mathbf P^4$}{P4} の新成分}

$A_0$ の代表は
\begin{align*}
 Q^{12}&=3xz/4,&Q^{13}&=0,&Q^{14}&=z^2+3x^2/4,\\
 Q^{23}&=-z^2/4,&Q^{24}&=x^2+2xz+3z^2-3xy/4+zw/4,&Q^{34}&=xz/4.
\end{align*}
$P=x_0^2J+Q$ を斉次化する。
全 $\mathcal Q$ では Jacobi 階数34、実軌道 $\operatorname{Sp}_4\times\mathbb G_m$ の次元10であり、
二次 jet の後、選択によらない三次障害が残る。
全 $\mathbf P^4$ では150係数の Jacobi 階数101、垂直方向階数24、
垂直方向・群作用・全体スカラーの合併階数48である。
従って射影接空間次元24に対し、実際の $\operatorname{PGL}_5$ 軌道次元は23である。
余剰方向を加えると階数49となり、三次障害の曖昧さを含む階数は101、
障害自身を加えると102となる。
ゆえに軌道閉包は全 $\mathbf P^4$ Poisson スキームの23次元成分で、一般横断環は
$K[\varepsilon]/(\varepsilon^3)$ である。

退化因子は
\[
 D=x_0F,\qquad F=x_0^4+x_0^2xz-z^4/4.
\]
$F$ は $x$ について原始一次式であり、
$\gcd(x_0^2z,x_0^4-z^4/4)=1$ だから既約である。
偏微分から $\operatorname{Sing}(F)=\{x_0=z=0\}$ が従い、
この唯一の特異平面の一般点での法線接錐は $x x_0^2z$、すなわち重複度型 $(2,1)$ である。
さらに $D\in I(\{x_0=z=0\})^4$ である。
平面に沿う4次消滅は射影 Grassmannian を用いた閉 incidence 条件であるため、
この軌道閉包は SNC 退化因子の開集合と交わらない。

v7の最小冪零例の四次因子
$F_{\min}=x_0^4+3x_0^2x^2+12x^3z$ の唯一の特異平面では、
一般法線接錐が $12z x^3$ という三重直線である。
従ってこの二つの23次元開軌道は異なる。
v6の開 $(2,2)$ 例は一般特異平面に沿う四次因子の重複度が2であり、今回の3と異なる。
他の既存成分との区別は、階数、次元、退化因子の既約分解、および上記閉 incidence 条件による。

\subsection{依存関係と未解決部分}

$G_1,G_2$ の名称・成分性・次元、および非冪零モジュラー開集合に交わる成分が
これらと八つの10次元成分で尽くされる事実には、v5の
\texttt{mixed\_modular\_theorem.txt} とそのv3/v4での半単純分類を用いる。
17固定ファイバー成分と11飽和層の網羅性・相互同一視にはv6の正確なファイバー分解を用いる。
今回の6新成分の全接空間・障害・実変形は、これらの固定ファイバー計算とは別に再構成した。

判別式 $\Delta=s^2-4\operatorname{Pf}(Q)$ は $G_1$ で $(Axy-Dzw)^2$、
$G_2$ で $A^2x^2y^2$ である。平方多項式の集合は閉集合なので、
これは両成分の\emph{被約台}についての閉条件である。
スキームの定義イデアルに平方条件が属する、あるいは冪零方向まで消去するという主張ではない。
$C_+,D_+$ の非平方判別式は追加の区別を与えるが、成分性と厚さは形式スライスで独立に証明した。

本節により $(2,2)$ 型の全11飽和層の $\mathcal Q$ 内での帰属は完了した。
全 $\mathbf P^4$ では、本節の $A_0$ に加え、$D_+$ も23次元・一般厚さ3の成分を与える
（定理\ref{thm:dplus-full}）。後者の全150係数の接空間と三次障害も別実装で監査済みである。
$C_+,C_-,U_1,AB$ の斉次化の全 $\mathbf P^4$ 内での成分帰属には、
さらに高次障害または実変形が必要である。
本節の結果に、最小冪零型と零モジュラー層の帰属の証明を合わせると、
全 $\mathcal Q$ の既約成分はちょうど22個となる。
全 $\mathbf P^4$ の既約成分の下界は55個であり、こちらの完全分類は未完である。
